\documentclass[10pt,twocolumn]{article}
\usepackage[T1]{fontenc}
\usepackage[utf8]{inputenc}
\usepackage{lmodern}
\usepackage[margin=1.8cm,top=2.4cm,bottom=2cm,columnsep=0.6cm]{geometry}
\usepackage{fancyhdr}
\usepackage{titlesec}
\usepackage{booktabs}
\usepackage{amsmath,amssymb,amsfonts}
\usepackage{microtype}
\usepackage{xcolor}
\usepackage{mdframed}
\usepackage{parskip}
\usepackage{enumitem}
\usepackage{hyperref}

\definecolor{rulered}{RGB}{180,30,30}
\definecolor{boxgray}{RGB}{245,245,245}
\definecolor{keywordgray}{RGB}{100,100,100}

\pagestyle{fancy}
\fancyhf{}
\fancyhead[L]{\textsc{\small Claude Mag}}
\fancyhead[C]{\small\textit{Machine Learning Research Digest}}
\fancyhead[R]{\small 2026-09-21}
\fancyfoot[C]{\small\thepage}
\renewcommand{\headrulewidth}{0.8pt}

\titleformat{\section}{\normalsize\bfseries\color{rulered}}{}{0pt}{}%
  [\vspace{-4pt}\color{rulered}\rule{\columnwidth}{0.4pt}]
\titlespacing{\section}{0pt}{8pt}{4pt}

\hypersetup{colorlinks=true,urlcolor=rulered,linkcolor=black}

\begin{document}
\setlength{\parindent}{0pt}
\setlength{\parskip}{4pt}

{\small\color{keywordgray}\textbf{Keywords:}
  agricultural AI \textbullet{}
  pixel-level grounding \textbullet{}
  multimodal large language model \textbullet{}
  plant disease recognition}\par\vspace{4pt}

{\LARGE\bfseries\raggedright
  AgriScope}\par\vspace{4pt}

{\large\itshape\raggedright
  A Pixel-Grounded Multimodal Model and 11-Million-Sample Dataset
  Bring Precision to Agricultural Image Understanding}\par\vspace{6pt}

{\small\textbf{By} Abderrahmene Boudiaf, Mohamad Alanssari, Irfan Hussain
  \& Sajid Javed (Khalifa University of Science and Technology)}\par\vspace{2pt}

{\small\color{keywordgray}arXiv:2609.20325 \textbullet{} ICML 2026}
\par\vspace{2pt}\noindent{\color{rulered}\rule{\columnwidth}{0.6pt}}\vspace{6pt}

Agricultural AI gets its first unified pixel-grounded multimodal model.
\textit{AgriScope} jointly supports image-level classification, region-level
referring segmentation, and pixel-level grounded caption generation within
a single framework trained on AgriGround---a new 500K-image,
11-million-sample dataset---lifting grounded caption CIDEr from 107.9 to
118.6 and mIoU from 55.0 to 59.4 over the best prior baseline on
the accompanying benchmark.

\begin{mdframed}[backgroundcolor=boxgray,linecolor=rulered,linewidth=1pt,
    innerleftmargin=6pt,innerrightmargin=6pt,
    innertopmargin=6pt,innerbottommargin=6pt]
\textbf{\small AT A GLANCE}\par\vspace{4pt}
\begin{description}[leftmargin=0pt,labelsep=4pt,itemsep=2pt,parsep=0pt,
    font=\small\bfseries]
  \item[Problem:] Agricultural image understanding requires fine-grained
    pixel-level localisation (disease lesions, pest bodies, growth
    stage features) that existing text-only MLLMs cannot provide.
  \item[Why it matters:] Actionable recommendations for precision
    agriculture (targeted spraying, selective harvesting) require knowing
    \emph{where} in the image the problem is, not just \emph{what} it is.
  \item[Method:] Biological-semantic encoder pre-trained on botanical
    taxonomies; multi-scale dense spatial features (FPN over BSE and
    DINOv2); mask transformer pixel decoder; AgriGround instruction-
    tuning dataset with 500K images and 11M samples.
  \item[Result:] GCG CIDEr 118.6 (+10.7), GCG mIoU 59.4 (+4.4),
    referring seg F1 73.2 (+5.4) over PixelLM baseline.
  \item[Evidence:] ICML 2026 evaluation across 6 agricultural task types;
    ablations over encoder, decoder, and dataset scale.
\end{description}
\end{mdframed}

\section{The Big Picture}

Precision agriculture aims to deliver the right intervention (pesticide,
fertiliser, pruning) to exactly the right location at exactly the right
growth stage. This requires a visual understanding system that can answer
not just ``does this plant have rust disease?'' but ``which leaves show
early-stage rust and where on each leaf are the uredinia?''

Multimodal large language models have advanced rapidly but their outputs
are purely textual: they name diseases and pests but do not localise them
in pixel coordinates. AgriScope closes this gap with a unified model that
takes an agricultural image (or image sequence) plus a natural-language
instruction and outputs both a text response \emph{and} pixel-level
segmentation masks for all referenced visual entities.

\section{Why Existing Approaches Fall Short}

\textbf{Text-only MLLMs} (GPT-4V, InternVL-2, LLaVA) produce written
diagnoses but cannot drive autonomous machinery that requires pixel masks
or bounding boxes. Their agricultural knowledge is also limited by
training on general web data rather than domain-curated imagery.

\textbf{Specialist agricultural classifiers} (PlantVillage CNNs,
domain-specific ViTs) excel at narrow tasks (identify this disease from
a leaf patch) but cannot answer open-ended agricultural questions,
handle novel species, or produce grounded multi-entity outputs.

\textbf{Grounding MLLMs} (LISA, PixelLM) support referring segmentation
for general-domain images but use standard ViT features that lack
biological specificity: they cannot distinguish powdery mildew from
leaf wax deposits, or aphid clusters from normal trichomes.

\section{Core Mechanism}

\textbf{Biological-Semantic Encoder (BSE).}
A ViT-L fine-tuned on 2.3M agricultural images annotated with
hierarchical biological labels (kingdom, class, order, family, genus,
species, disease type, disease stage). The BSE produces feature tokens
$\{f_i^{\text{bio}}\}$ embedding biological category memberships
alongside spatial features.

\textbf{Dense Spatial Representations (DSR).}
A multi-scale feature pyramid fuses BSE tokens with a frozen DINOv2
encoder:
\[
F_{\text{DSR}} = \text{FPN}(f^{\text{bio}}, f^{\text{dino}})
               \in \mathbb{R}^{C \times H/4 \times W/4},
\]
providing spatially dense features at three resolutions for both
coarse recognition and fine lesion localisation.

\textbf{Pixel Decoder.}
A mask transformer decoder takes query embeddings $\{q_k\}$ derived
from the LLM output's \texttt{[SEG]} tokens and predicts binary masks:
\[
M_k = \text{PD}(q_k, F_{\text{DSR}}) \in \{0,1\}^{H \times W}.
\]
When the LLM generates \texttt{[SEG]}, the last hidden state at that
position feeds directly into the pixel decoder as $q_k$.

\section{Technical Formulation}

Training minimises:
\[
\mathcal{L} = \mathcal{L}_{\text{LM}}
            + \lambda_{\text{BCE}}\mathcal{L}_{\text{BCE}}
            + \lambda_{\text{dice}}\mathcal{L}_{\text{dice}},
\]
with $\lambda_{\text{BCE}}=1.0$, $\lambda_{\text{dice}}=0.5$.

Binary cross-entropy on predicted mask $\hat{M}_k$ vs.\ ground-truth $M_k^*$:
\[
\mathcal{L}_{\text{BCE}} = -\sum_p
  \bigl[M_k^*(p)\log\hat{M}_k(p)
  + (1{-}M_k^*(p))\log(1{-}\hat{M}_k(p))\bigr].
\]
Dice loss for shape overlap:
\[
\mathcal{L}_{\text{dice}} = 1 -
  \frac{2\sum_p M_k^*(p)\hat{M}_k(p)}
       {\sum_p M_k^*(p) + \sum_p \hat{M}_k(p)}.
\]

\section{Learning or Inference Procedure}

\begin{enumerate}[itemsep=2pt,parsep=0pt]
  \item Pre-train BSE on 2.3M agricultural images with hierarchical
    biological classification (warm-start from ViT-L ImageNet checkpoint).
  \item Train AgriScope on AgriGround for 3 epochs with frozen DINOv2
    and trainable BSE, LLM adapters (LoRA), and pixel decoder.
  \item Optional task-specific fine-tuning for 1 epoch on narrow
    subsets (e.g., disease detection only).
  \item Inference: input image + instruction; LLM generates text
    interleaved with \texttt{[SEG]} tokens; pixel decoder predicts
    a mask for each \texttt{[SEG]}.
\end{enumerate}

AgriGround construction: 500K source images from PlantVillage,
iNaturalist-Agriculture, FGVC-Plant, and CropDiseaseDB; captions
generated by GPT-4V with botanical system prompts; grounded with
Grounding-DINO\,+\,SAM2; 11M instruction-following samples across
6 task types (disease, pest, species, growth, weed, nutritional deficiency).

\section{What the Guarantee Says}

No formal guarantees. The ICML reviewers cited the breadth of the
AgriGround evaluation: 6 task types, 3 cross-dataset generalisation
settings, and both in-distribution and zero-shot conditions, providing
strong empirical evidence that improvements are not evaluation artefacts.

\section{Experimental Findings}

Main benchmark: AgriGround-Test (10K held-out samples).

\begin{center}
\small
\begin{tabular}{lcccc}
\toprule
Method & CIDEr & mIoU & Seg F1 & VQA Acc \\
\midrule
LLaVA-1.5 & 82.3 & 43.7 & --- & 71.4 \\
InternVL-2 & 97.1 & 49.8 & --- & 78.9 \\
LISA & 101.4 & 52.3 & 64.1 & 74.2 \\
PixelLM & 107.9 & 55.0 & 67.8 & 76.5 \\
\textbf{AgriScope} & \textbf{118.6} & \textbf{59.4} & \textbf{73.2} & \textbf{82.7} \\
\bottomrule
\end{tabular}
\end{center}

CIDEr and mIoU for grounded caption generation (GCG). LLaVA and
InternVL-2 do not support pixel-level output (Seg F1 not applicable).
AgriScope is the first model to achieve over 59 mIoU on this benchmark.

\section{Ablations and Interpretation}

\textbf{Standard ViT encoder (no BSE):} CIDEr $-6.8$; the BSE's
biological specificity is critical for fine-grained disease staging.

\textbf{Single-resolution features (no FPN):} mIoU $-3.1$; multi-scale
features are needed to localise both large leaf lesions and small pest bodies.

\textbf{Without AgriGround (existing datasets only):} CIDEr $-14.3$;
breadth and scale of agricultural supervision are the primary driver.

\textbf{Dataset scaling:} 100K samples $\to$ 250K: $+6.2$ CIDEr;
250K $\to$ 500K: $+2.1$ CIDEr. Diminishing returns but still positive,
suggesting continued benefit from data collection.

\textbf{Decoder variant:} Mask transformer vs.\ FPN upsampling:
$+2.7$ mIoU; mask transformer's cross-attention to text queries is
essential for language-conditioned segmentation.

\section*{Reference}

Abderrahmene Boudiaf, Mohamad Alanssari, Irfan Hussain, Sajid Javed.
\textbf{AgriScope: Pixel-Grounded Multimodal Understanding for Agricultural
Images.}
arXiv:2609.20325, September 2026. ICML 2026.
\url{https://arxiv.org/abs/2609.20325}

\end{document}
